\subsection{Conditional measures of PfPPs}
Let $X$ and $\pr{X}$ be disjoint finite subsets in $\mfrak{X}$ and take a cylinder set
\begin{equation*}
	C(X,\pr{X})=\{\omega\in \Omega\,|\, X\subset \omega,\, \pr{X}\cap\omega=\varnothing\},
\end{equation*}
which consists of configurations such that every points in $X$ are occupied and those in $\pr{X}$ are unoccupied.
We identify the cylinder set $C(X,\pr{X})$ with $\Omega (\mfrak{X}\backslash (X\sqcup\pr{X})):=\{0,1\}^{\mfrak{X}\backslash (X\sqcup\pr{X})}$ via the map
\begin{equation*}
	F_{X,\pr{X}}\colon \ C(X,\pr{X})\to \Omega (\mfrak{X}\backslash (X\sqcup\pr{X}));\qquad \omega\mapsto \omega\backslash X.
\end{equation*}
For a probability measure $M$ on $(\Omega,\Sigma)$, assume that the cylinder set has strictly positive weight: $M(C(X,\pr{X}))>0$.
We define the conditional measure of $M$ on $(X,\pr{X})$ by
\begin{equation*}
	M_{X,\pr{X}}:=(F_{X,\pr{X}})_{\ast}\frac{M\big|_{C(X,\pr{X})}}{M(C(X,\pr{X}))}.
\end{equation*}

We focus on PfPPs associated with projection operators.
For a subset $A\subset\mfrak{X}$, we set
\begin{equation*}
	\mcal{K}_{A}^{+}=\overline{\mrm{Span}\{(e_{x},0)\,|\,x\in A\}},\qquad \mcal{K}^{-}_{A}=\overline{\mrm{Span}\{(0,e_{x})\,|\,x\in A\}}
\end{equation*}
and $\mcal{K}_{A}=\mcal{K}_{A}^{+}\oplus\mcal{K}_{A}^{-}$, which is the direct sum of Hilbert spaces.
It is obvious from the definition that $\Gamma$ preserves the subspace $\mcal{K}_{A}$. Thus, we can write $\Gamma_{A}:=\Gamma|_{\mcal{K}_{A}}$.

Let $X$ and $\pr{X}$ be finite disjoint subsets of $\mfrak{X}$ and take a projection operator $P\in \mrm{Gr}(\mcal{K},\Gamma)$.
Then, $P\mcal{K}\subset\mcal{K}$ is a closed subspace. We define a new closed subspace
\begin{equation*}
	P\mcal{K}_{X,\pr{X}}:=\big(P\mcal{K}+\big(\mcal{K}_{X}^{+}\oplus \mcal{K}^{-}_{\pr{X}}\big)\big)\cap \mcal{K}_{\mfrak{X}\backslash (X\sqcup\pr{X})}
\end{equation*}
in $\mcal{K}_{\mfrak{X}\backslash (X\sqcup\pr{X})}$ and a projection operator $P_{X,\pr{X}}$ as the orthogonal projection onto $P\mcal{K}_{X,\pr{X}}$ in $\mcal{K}_{\mfrak{X}\backslash (X\sqcup\pr{X})}$.

\begin{Lemma}\label{lem:conditional_projection}
Let $X$ and $\pr{X}$ be finite disjoint subsets of $\mfrak{X}$. Then, we have
\begin{equation*}
	P_{X,\pr{X}}\in\mrm{Gr}\big(\mcal{K}_{\mfrak{X}\backslash(X\sqcup\pr{X})},\Gamma_{\mfrak{X}\backslash (X\sqcup\pr{X})}\big).
\end{equation*}
\end{Lemma}

 Let us also introduce the notion of regularity.
 \begin{Definition} Let $X$ and $\pr{X}$ be finite disjoint subsets of $\mfrak{X}$.
 A projection operator $P\in \mrm{Gr}(\mcal{K},\Gamma)$ is said to be $(X,\pr{X})$-regular if
 \begin{equation*}
 	P\mcal{K}\cap \big(\mcal{K}_{X}^{+}\oplus \mcal{K}_{\pr{X}}^{-}\big)=\{0\}.
 \end{equation*}
 \end{Definition}

 The following result is an analogue of \cite[Proposition~6.13]{Olshanski2020} for PfPPs.
 \begin{Theorem}\label{thm:conditional_projection}
 Let $X$ and $\pr{X}$ be finite disjoint subsets of $\mfrak{X}$.
 Assume that a projection operator $P\in\mrm{Gr}(\mcal{K},\Gamma)$ is $(X,\pr{X})$-regular.
 Then, we have
 \begin{equation*}
 	M^{P}_{X,\pr{X}}:=\big(M^{P}\big)_{X,\pr{X}}=M^{P_{X,\pr{X}}}.
 \end{equation*}
 In particular, it is a PfPP associated with a projection operator.
 \end{Theorem}

 \begin{Remark}
A similar problem has been studied in~\cite{BufetovCundenQiu2019}. Our result describes the reduction of a projection operator, which is new.
 \end{Remark}
